\documentclass[a4paper,11pt]{ltjsarticle}
\usepackage[haranoaji]{luatexja-preset}
\usepackage{amsmath,amssymb}
\usepackage[top=12mm,bottom=10mm,left=14mm,right=14mm]{geometry}
\usepackage{array}
\usepackage{tikz}
\usetikzlibrary{calc}
\pagestyle{empty}
\setlength{\parindent}{0pt}
\newlength{\qcol}\setlength{\qcol}{0.47\textwidth}
\newlength{\qgap}
\newlength{\anscolw}\setlength{\anscolw}{0.30\textwidth}
\newlength{\ansh}
\newcommand{\Q}[2]{\parbox[t]{\qcol}{\raggedright (#1)\hspace{0.6em}$\displaystyle #2$}}
\newcommand{\QR}[4]{\Q{#1}{#2}\hfill\Q{#3}{#4}\par\vspace{\qgap}}
\newcommand{\anscell}[1]{\rule[-\ansdrop]{0pt}{\ansh}(#1)}
\newlength{\ansdrop}

\begin{document}
{\large\bfseries 数学テスト　11}\hfill 名前(\underline{\hspace{32mm}})\par\vspace{3mm}
■（1）〜（12）の計算をしなさい。（13）、（14）は連立方程式を解きなさい。\par\vspace{3mm}
\setlength{\qgap}{11mm}
\QR{1}{\dfrac{3}{4}(8x-16y)}{2}{16a^3b\div2ab\div4a}
\QR{3}{(5x-3y)+(-4x+y)}{4}{(5a^2+2a-3)-(-2a^2-4a+1)}
\QR{5}{12a^2\div4a\times3ab}{6}{(28x-16y)\div4}
\QR{7}{\dfrac{a-3b}{2}-\dfrac{3a+b}{3}}{8}{9a+3-\{5b+(2a-4b)-1\}}
\QR{9}{-32a^2b^2\div8ab}{10}{(3x+2y)-(5x-4y)+(x-6y)}
\QR{11}{(2x)^2\times xy}{12}{2(x+2y)+3(2x-3y)}
\QR{13}{\left\{\begin{array}{l} 5x+4y=6 \\ x-3y=5 \end{array}\right.}{14}{\left\{\begin{array}{l} 0.6x+0.5y=2 \\ 2x+3y=4 \end{array}\right.}
\vfill
\setlength{\ansh}{12mm}\setlength{\ansdrop}{8mm}%
\begin{center}\begin{tabular}{|p{\anscolw}|p{\anscolw}|p{\anscolw}|}\hline
\anscell{1} & \anscell{2} & \anscell{3} \\\hline
\anscell{4} & \anscell{5} & \anscell{6} \\\hline
\anscell{7} & \anscell{8} & \anscell{9} \\\hline
\anscell{10} & \anscell{11} & \anscell{12} \\\hline
\anscell{13} & \anscell{14} &  \\\hline
\end{tabular}\end{center}

\end{document}
